\chapter{Lecture 07}

\hfill {\footnotesize \textit{Date: 04/03/26}}
\section{Generation of Modules}
Let \(M\) be a left  \(R\)-module and  \(m \in M\). Then 
 \[
\langle m \rangle \coloneq Rm = \{rm  \mid r \in R\} 
\] is a submodule of \(M\) \emph{generated by}  \(m\).

Suppose \(\{x_\alpha\}_{\alpha \in \mathcal{A}}\) is an arbitrary subset of \(M\), and that \(M = \sum_{\alpha \in \mathcal{A}} R x_\alpha\). Then  \(M\) is said to be \emph{generated by the set} \(\{x_\alpha\}_{\alpha \in \mathcal{A}}\). An \(R\)-module is \emph{finitely generated} if it generated by a finite set.

\section{Direct Product \& Direct Sum}
\begin{definition}[Direct Product]
     Suppose \(\{M_\alpha\}_{\alpha \in \mathcal{A}}\) is an arbitrary collection of \(R\)-modules. Then their \emph{direct product} is the \(R\)-module 
    \[
        \prod_{\alpha \in \mathcal{A}} M_\alpha  \coloneq \left\{(m_\alpha)_{\alpha \in \mathcal{A}}  \mid m_\alpha \in M_\alpha \text{ for all } \alpha \in \mathcal{A}\right\}
    \]
    equipped with component-wise addition and scalar multiplication.
\end{definition}
\begin{definition}[Direct Sum]
    Suppose \(\{M_\alpha\}_{\alpha \in \mathcal{A}}\) is an arbitrary collection of \(R\)-modules. Then their \emph{direct sum} is the \(R\)-module 
    \[
         \bigoplus_{\alpha \in \mathcal{A}} M_\alpha  \coloneq \left\{(m_\alpha)_{\alpha \in \mathcal{A}}  \mid m_\alpha \in M_\alpha\text{ and all but finitely many } m_\alpha's \text{ are zero}\right\}
    \]
    equipped with component-wise addition and scalar multiplication.
\end{definition}
\begin{remark}
    If the indexing set \(\mathcal{A}\) is finite, then direct product and direct sum coincide.
\end{remark}
\section{Freely Generated Modules}
\begin{definition}
    An \(R\)-module \(M\) is free if 
     \[
         M \cong \bigoplus_{\alpha \in \mathcal{A}} M_\alpha, \quad \text{where each } M_\alpha \cong R \text{ as } R\text{-modules}.
    \] 
\end{definition}
\begin{remark}
    A finitely generated free module is isomorphic to \(R \oplus \cdots \oplus R = R^n\).
\end{remark}
%\begin{proof}
%    It is immediate that a free module M module is isomorphic to \(\bigoplus_{\alpha \in \mathcal{A}} R\) for some indexing set \(\mathcal{A}\). We have to show that \(\mathcal{A}\) is finite.

%    \(\bigoplus_{\alpha \in \mathcal{A}} R\) has a natural basis \(\mathcal{E} = \{e_\beta\}_{\beta \in \mathcal{A}} \) where each \(e_\beta = (x_\alpha)_{\alpha \in \mathcal{A}}\), where \(x_\alpha = 1\) if  \(\alpha = \beta\), and \(x_\alpha = 0\) otherwise.

%    Since \(M\) is finitely generated, so is \(\bigoplus_{\alpha \in \mathcal{A}} R\), i.e., generated by a finite set \(\{(x_\alpha)_{\alpha \in \mathcal{A}}^{(i)}\}_{i=1}^n\), where \((x_\alpha)_{\alpha \in \mathcal{A}}^{(i)} \in \bigoplus_{\alpha \in \mathcal{A}} R\). 

%    Then \((x_\alpha)_{\alpha \in \mathcal{A}}^{(i)} = \sum_{\alpha \in \mathcal{A}} r_\alpha^{(i)} e_\alpha\) where \(r_\alpha^{(i)} = 0\) for all but finitely many  \(\alpha \in \mathcal{A}\). Define
%    \[
%        \mathcal{F} \coloneq \{e_\alpha \mid r_\alpha^{(i)} \neq 0 \text{ in the sum }  (x_\alpha)_{\alpha \in \mathcal{A}}^{(i)} = \sum_{\alpha \in \mathcal{A}} r_\alpha^{(i)} e_\alpha \text{ for some } i = 1, \ldots, n\}. 
%    \] 
%    Since \((x_\alpha)_{\alpha \in \mathcal{A}}^{(i)} = \sum_{\alpha \in \mathcal{A}} r_\alpha^{(i)} e_\alpha\) is finite,  \(\mathcal{F}\) is finite as well. We have \(\mathcal{F} = \mathcal{E}\) trivially. Since \(\mathcal{E}\) is indexed by \(\mathcal{A}\), so \(\mathcal{E}\) being finite implies \(\mathcal{A}\) is finite as well, completing the proof.
%\end{proof}
\begin{proposition}
    \(M\) is finitely generated  \(R\)-module if and only if  \(M\) is isomorphic to a quotient of  \(R^n\) for some  integer \(n > 0\).
\end{proposition}
\begin{proof}
(\(\implies\)) Suppose \(M\) is finitely generated. Then there exists \(x_1, \ldots, x_n \in M\) such that 
\[
    m \in M \implies m = \sum_{i=1}^{n} r_i x_i; \quad r_i \in R.
\] 
Define
\[
    \Phi: \bigoplus_{i=1}^n R \longrightarrow M 
\]
\[
    (r_1, \ldots, r_n) \mapsto \sum_{i=1}^{n} r_i x_i.
\] 
It is immediate that \(\Phi\) is an  \(R\)-module homomorphism and \(\Phi\) is surjective. 

Using first isomorphism theorem, we have
\[
    \left(\bigoplus_{i-1}^n R\right) \Biggm/ (\ker \Phi) \cong M.
\] 
(\(\Longleftarrow\)) Suppose \(M \cong \left(\bigoplus_{i=1}^n R\right)/N\) for some submodule \(N\) via an isomorphism  \(\Phi\). 
% https://q.uiver.app/#q=WzAsMyxbMCwwLCJcXGJpZ29wbHVzX3tpPTF9Xm5SIl0sWzEsMCwiXFxsZWZ0KFxcYmlnb3BsdXNfe2k9MX1ebiBSXFxyaWdodCkvTiJdLFsyLDAsIk0iXSxbMCwxLCJcXHBpIl0sWzEsMiwiXFxQaGkiXSxbMCwyLCIiLDIseyJjdXJ2ZSI6M31dXQ==
\[\begin{tikzcd}
	{\bigoplus_{i=1}^nR} & {\left(\bigoplus_{i=1}^n R\right)/N} & M
	\arrow["\pi", from=1-1, to=1-2]
\arrow["{\Phi \circ \pi}", bend right, from=1-1, to=1-3]
	\arrow["\Phi", from=1-2, to=1-3]
\end{tikzcd}\]
\(\{e_i\}_{i=1}^n\) is a generating set for \(\bigoplus_{i=1}^n R\), so  \(\{(\Phi \circ \pi)(e_i)\}\) is a generating set of \(M\) since \(\Phi \circ \pi\) is surjective. 
\end{proof}
\begin{proposition}\label{prop:fingeneratedmodpoly}
    Suppose \(M\) is finitely generated \(R\)-module, and let \(\mathfrak{a}\) be an ideal of \(R\). Let \(\phi\) be an \(R\)-module endomorphism such that \(\phi(M) \subseteq \mathfrak{a} M\). Then \(\phi\) satisies an equation of the form \[\phi^n + \phi^{n-1} a_1 + \cdots + a_n = 0\]
    where \(a_1, \ldots, a_n \in \mathfrak{a}\). 
\end{proposition}
\begin{proof}
    Let \(x_1, \ldots, x_n\) be generators of \(M\). Then each \(\phi(x_i) \in \mathfrak{a}M\), so that we have \(\phi(x_i) = \sum_{j=1}^n a_{ij}x_j (1 \le i \le n; a_{ij} \in \mathfrak{a})\), i.e., \[\sum_{j=1}^n (\delta_{ij}\phi - a_{ij})x_j = 0\]
    where \(\delta_{ij}\) is Kronecker delta function, and that gives us the following matrix equation:
    \[[\delta_{ij} \phi - a_{ij}] \begin{pmatrix}
    x_1\\
    \vdots\\
    x_n
\end{pmatrix} = 0.\]
Multiplying adjoint matrix of \([\delta_{ij}\phi - a_{ij}]\) both sides, we have
\begin{align*}
    &\text{Adj}[\delta_{ij} \phi - a_{ij}] [\delta_{ij} \phi - a_{ij}] \begin{pmatrix}
    x_1\\
    \vdots\\
    x_n
\end{pmatrix}= 0\\
    \implies &\det\left([\delta_{ij} \phi - a_{ij}]\right) [\delta_{ij} \phi - a_{ij}]^{-1} [\delta_{ij} \phi - a_{ij}]\begin{pmatrix}
    x_1\\
    \vdots\\
    x_n
\end{pmatrix}  = 0\\
&\text{(because for all $n \times n$ invertible matrices $M, \det(M) M^{-1} = \text{Adj}(M)$)}\\ 
        \implies &\det([\delta_{ij} \phi - a_{ij}]) \begin{pmatrix}
    x_1\\
    \vdots\\
    x_n
\end{pmatrix} = 0.
\end{align*}
So, it follows that \(\det([\delta_{ij}\phi - a_{ij}])\) annihilates each \(x_i\), hence is the zero endomorphism of \(M\). Expanding out \(\det([\delta_{ij}\phi - a_{ij}])\) gives us an equation of the required form. 
\end{proof}
